\section{Zadanie 10}
\autor{Filip Manijak}
Niech \((\Omega, \Sigma)\) i \((\Omega', \Sigma')\) będą przestrzeniami mierzalnymi oraz \(\sigma(\mathcal{G}) = \Sigma'\) dla pewnego \(\mathcal{G} \subseteq \mathcal{P}(\Omega')\) (innymi słowy \(\Sigma'\) jest generowana przez \(\mathcal{G}\)). Udowodnij, że \(f : \Omega \to \Omega'\) jest mierzalna wtedy i tylko wtedy, gdy \(f^{-1}(G) \in \Sigma\) dla każdego \(G \in \mathcal{G}\).

\subsection*{Rozwiązanie}
Sprawadza się to, do udowodnienia że:
\[f^{-1}(G) \in \Sigma \iff f^{-1}(\sigma(G)) \in \Sigma\]
\subsection{\(\Leftarrow\)}
Triv, bo \(G \subset \sigma(G)\)
\subsection{\(\implies\)}
Tu plan jest taki, że pokażemy że wszystkie zbiory o mierzalnym przeciwobrazie są \(\sigma\)-algebrą (oznaczmy ją \(M\) jak mierzalne). Wiemy jednak, że \(G\) należy do tej sigma algebry, a to oznacza że \(\sigma(G) \subset M\) bo \(\sigma(G)\) to przecięcie \(M\) z czymś tam (nie interesuje nas czym). No to teraz ta trudniejsza cześć, pokazanie tego
\subsection{Dowód że \(\{X: f^{-1}(X) \in \Sigma\} = M\) jest \(\sigma\)-algebrą}
Sprawdzamy własności:
\begin{itemize}
	\item \(\Omega' \in M\), tak ponieważ \(f^{-1}(\Omega') = \Omega \in \Sigma\)
	\item \(A \in M \implies \Omega'/A \in M\) też, ponieważ \(f^{-1}(\Omega'/A) = f^{-1}(\Omega')/f^{-1}(A)\) co już jest w \(M\) bo \(\Sigma\) to algebra
	\item Suma przeciwobrazów podobnie się przenosi, więc ten sam argument.
	\item Suma przeliczalnej ilości zachowuje się trywialnie tak samo.
\end{itemize}
